\paragraph{Bit-Flip Errors}\label{sec:bit-flip-errors}

While dephasing affects the \textbf{relative phase} of a quantum state, amplitude dampig affects the \textbf{probabilities (populations)} of the basis states. In practice, \definition{Amplitude Damping} occurs when the \textbf{qubit exchanges energy with the environment}. For example, initially the state:
\begin{equation*}
    \ket{\psi} = \frac{\ket{0} + \ket{1}}{\sqrt{2}} = \ket{+}
\end{equation*}
With probabilities:
\begin{equation*}
    P(0) = P(1) = \frac{1}{2}
\end{equation*}
After amplitude damping, the state might evolve into:
\begin{equation*}
    \ket{\psi} = \sqrt{0.8}\ket{0} + \sqrt{0.2}\ket{1}
\end{equation*}
With probabilities:
\begin{equation*}
    P(0) = 0.8 \qquad P(1) = 0.2
\end{equation*}
The population of \ket{1} has decreased. This is amplitude damping. Specifically, this \textbf{form of amplitude damping} is called \definitionWithSpecificIndex{Energy Relaxation}{Relaxation}{Energy Relaxation} because the qubit loses energy to the environment. In this case, the state tends to relax towards the ground state \ket{0}.

\highspace
\textcolor{Green3}{\faIcon{question-circle} \textbf{Why does Energy Relaxation occur?}} Energy relaxation occurs because the qubit is not perfectly isolated from its surroundings. The qubit can interact with the environment, exchanging energy in the process. For example, if the qubit is implemented using a superconducting circuit, it can lose energy by emitting photons into the surrounding electromagnetic field. This process causes the excited state \ket{1} to decay towards the ground state \ket{0}, leading to amplitude damping.

\highspace
\textcolor{Red2}{\faIcon{exclamation-triangle} \textbf{Another type of amplitude damping: Thermal Excitation.}} In some cases, the \textbf{qubit can also gain energy from the environment}, leading to \definitionWithSpecificIndex{Thermal Excitation}{Excitation}{Thermal Excitation}. For example, if the \hl{qubit is in a warm environment} (i.e., at a higher temperature), it can \hl{absorb thermal energy} and transition from the ground state \ket{0} to the excited state \ket{1}. This process also \hl{leads to amplitude damping because it changes the population of the basis states.}

\highspace
\textcolor{Green3}{\faIcon{question-circle} \textbf{Where does the Bit-Flip come from?}} Suppose Energy Relaxation occurs:
\begin{equation*}
    \ket{1} \to \ket{0}
\end{equation*}
What happened is that the logical value changed from $1$ to $0$. Instead, suppose Thermal Excitation occurs:
\begin{equation*}
    \ket{0} \to \ket{1}
\end{equation*}
Now the logical value changed from $0$ to $1$. In both cases, the logical value of the qubit has flipped. This is exactly what the Pauli-$X$ gate does:
\begin{equation*}
    X = \begin{bmatrix}
        0 & 1 \\
        1 & 0
    \end{bmatrix}
\end{equation*}
Which acts as follows:
\begin{equation*}
    X\ket{0} = \ket{1} \qquad X\ket{1} = \ket{0}
\end{equation*}
Therefore, we can model amplitude damping as a \definition{Bit-Flip Error} because it changes the logical value of the qubit. Just like with phase-flip errors, a single application of the $X$ gate does not necessarily destroy the quantum state, but in a real quantum computer, interactions with the environment can introduce many random bit-flip errors over time, leading to significant amplitude damping and loss of quantum information.

\highspace
\textcolor{DarkOrange3}{\faIcon{exclamation-triangle} \textbf{Important Caveat.}} Amplitude damping is not literally a bit-flip; we use \hl{bit-flips as a convenient mathematical model}. A Pauli-$X$ gate is a perfect bit-flip and treats the two states symmetrically $\ket{0} \iff \ket{1}$. Same behavior in both directions. Instead, the nature of amplitude damping is asymmetric. Usually, the qubit \ket{1} tends to decay towards \ket{0} (energy relaxation), and the reverse process (thermal excitation) is less common. Because \ket{1} is typically the \textbf{higher-energy state}, while \ket{0} is the \textbf{lower-energy state}, and the qubit naturally tends to lose energy to the environment rather than gain it. Therefore, while we can model amplitude damping as a bit-flip error for simplicity, it's important to remember that the underlying physical processes are more complex and not perfectly symmetric.

\textcolor{DarkOrange3}{\faIcon{exclamation-triangle} \textbf{Important Caveat.}} Amplitude damping is not literally a bit-flip; rather, we use \hl{bit-flips as a convenient mathematical model}. A Pauli-$X$ gate is a perfect bit-flip and treats the two states symmetrically:
\begin{equation*}
    \ket{0} \leftrightarrow \ket{1}
\end{equation*}
In other words, the behavior is identical in both directions. Amplitude damping, however, is generally asymmetric. Usually, the qubit $\ket{1}$ tends to decay towards $\ket{0}$ (energy relaxation), while the reverse process $\ket{0}\rightarrow\ket{1}$ (thermal excitation) is less likely. This happens because $\ket{1}$ is typically the \textbf{higher-energy state}, whereas $\ket{0}$ is the \textbf{lower-energy state}. Consequently, a qubit naturally tends to lose energy to the environment rather than gain it. Therefore, although amplitude damping is often modeled and corrected as an $X$-type (bit-flip) error, the underlying physical mechanism is more complex and not perfectly symmetric. In summary, a bit-flip describes \emph{what happens} to the logical value of the qubit, while amplitude damping describes \emph{why it happens} physically.